\chapter{繼承與多型}
\label{ch:inheritance}

第~5、6~章的類別各自獨立：\texttt{Account} 是帳戶，\texttt{Fraction} 是
分數，彼此沒有關係。 實際的程式常有一群「同一種東西的不同版本」：圓和矩
形都是形狀，汽車和機車都是車。 它們有共同的部分（名字、都能算面積），也有
各自不同的部分（半徑、寬和高、算面積的公式）。

這一章先用\textbf{繼承}（inheritance）表達「圓是一種形狀」，看衍生類別
怎麼建立和銷毀。 接著用\textbf{虛擬函式}讓同一行 \texttt{s.area()} 依物件
實際的型別呼叫不同的函式，這叫\textbf{多型}（polymorphism）。 後半章講抽
象類別、虛擬解構子，以及怎麼把不同型別的物件放進同一個 \texttt{vector}。
最後一節講什麼時候不該用繼承。

\section{衍生類別}
\label{sec:derived}

\begin{crecap}[型別標記與共同開頭]
C 要表達「圓和矩形都是形狀」，常見的寫法是在結構裡放一個\textbf{型別標
記}，用 \texttt{switch} 分開處理：

\begin{code}{c}{用型別標記區分形狀}
enum Kind { CIRCLE, RECT };

struct Shape {
    enum Kind kind;        /* 型別標記 */
    const char *name;
    double a, b;           /* 圓：a 是半徑；矩形：a、b 是寬和高 */
};

double area(const struct Shape *s) {
    switch (s->kind) {
    case CIRCLE: return 3.14159 * s->a * s->a;
    case RECT:   return s->a * s->b;
    }
    return 0;
}
\end{code}

\texttt{a}、\texttt{b} 代表什麼，只能靠註解約定。 多一種形狀，就要找出每
一個對 \texttt{kind} 做 \texttt{switch} 的函式，一個一個補上。

另一種寫法是把共同的部分做成一個結構，放在每種形狀的第一個成員，再把「算
面積」做成函式指標：

\begin{code}{c}{共同開頭加函式指標}
struct Shape {
    const char *name;
    double (*area)(const struct Shape *self);
};

struct Circle {
    struct Shape base;     /* 共同的部分放在第一個成員 */
    double r;
};

/* c 是一個 struct Circle，建立時把 c.base.area 設成 circle_area */
struct Shape *p = (struct Shape *)&c;   /* 第一個成員的位址就是 c 的位址 */
printf("%s %g\n", p->name, p->area(p)); /* 圓 3.14159 */
\end{code}

\texttt{circle\_area} 收 \texttt{struct Shape *}，自己轉回
\texttt{struct Circle *} 再算。 拿到 \texttt{struct Shape *} 的程式不用知
道它是哪一種形狀，呼叫 \texttt{p->area(p)} 就會執行對應的函式。 代價是到
處都要強制轉型，函式指標也要自己填，填錯了編譯器不會發現。
\end{crecap}

C++ 的繼承做的是第二種寫法的前半：把共同的部分寫成一個類別，其他類別「繼
承」它。 被繼承的叫\textbf{基底類別}（base class），繼承的叫\textbf{衍生
類別}（derived class）。 函式指標那一半是第~\ref{sec:virtual}~節的虛擬
函式。

\begin{example}[label=ex:derived]{cpp}{(derived.cpp) Circle 和 Rect 都繼承 Shape}
#include <iostream>
#include <string>
using namespace std;

class Shape {
public:
    Shape(const string &name) : name(name) {}
    const string &getName() const { return name; }

private:
    string name;
};

class Circle : public Shape {
public:
    Circle(double r) : Shape("圓"), r(r) {}
    double area() const { return 3.14159 * r * r; }

private:
    double r;
};

class Rect : public Shape {
public:
    Rect(double w, double h) : Shape("矩形"), w(w), h(h) {}
    double area() const { return w * h; }

private:
    double w, h;
};

int main() {
    Circle c(1);
    Rect r(3, 4);
    cout << c.getName() << ' ' << c.area() << '\n';
    cout << r.getName() << ' ' << r.area() << '\n';

    const Shape &s = r;        // Rect 也是一個 Shape
    cout << s.getName() << '\n';
    // s.area();               // 編譯失敗：Shape 沒有 area
    return 0;
}
\end{example}

\iosample{}{
圓 3.14159 \\
矩形 12 \\
矩形 \\
}

\begin{itemize}
\item \texttt{class Circle : public Shape} 表示 \texttt{Circle} 繼承
\texttt{Shape}。 每個 \texttt{Circle} 物件裡都有一份 \texttt{Shape}（跟
C 的 \texttt{base} 成員一樣），再加上自己的 \texttt{r}。
\texttt{Shape} 的 \texttt{public} 成員函式也變成 \texttt{Circle} 的，所以
可以寫 \texttt{c.getName()}。
\item 建構子不會繼承。 \texttt{Circle} 的建構子在初始化列表寫
\texttt{Shape("圓")}，用 \texttt{Shape} 的建構子建立裡面那一份
\texttt{Shape}，再建立自己的成員 \texttt{r}。
\item \texttt{name} 是 \texttt{Shape} 的 \texttt{private} 成員。
\texttt{Circle} 的物件裡有它，但 \texttt{Circle} 的成員函式不能直接用，只
能透過 \texttt{getName()}。 \texttt{private} 對衍生類別也一樣有效，
\texttt{Shape} 的不變量還是只要檢查 \texttt{Shape} 自己的函式。
\item \texttt{const Shape \&s = r;} 把 \texttt{Rect} 綁到 \texttt{Shape}
的參考上，同樣可以寫 \texttt{Shape *p = \&r;}。 這就是 C 回顧框裡的
\texttt{(struct Shape *)\&c}，只是不用轉型，編譯器知道
\texttt{Rect} 是一種 \texttt{Shape}。 反過來不行：\texttt{Shape} 不一定是
\texttt{Rect}。
\item 透過 \texttt{s} 只能用 \texttt{Shape} 有的成員。 \texttt{s.area()}
報 \texttt{'const class Shape' has no member named 'area'}。 讓
\texttt{s.area()} 能用、而且呼叫到 \texttt{Rect} 的 \texttt{area}，是
第~\ref{sec:virtual}~節的內容。
\end{itemize}

\texttt{public} 繼承的意思是「\texttt{Circle} 是一種 \texttt{Shape}」
（is-a）：任何需要 \texttt{Shape} 的地方，都可以交給它一個
\texttt{Circle}。 決定要不要用繼承時，先問這句話說不說得通。
第~\ref{sec:composition}~節會再回來談這件事。

\subsection{\texttt{protected}}

衍生類別有時候需要碰基底類別的內部，但又不想開放給所有人。 這時用第三種
存取標籤 \texttt{protected}（保護）：\texttt{protected} 成員對外面的程式跟
\texttt{private} 一樣，衍生類別的成員函式卻可以用。 例如在 \texttt{Shape}
加一段 \texttt{protected: void setName(const string \&n) \{ name = n; \}}，
\texttt{Circle} 的成員函式就可以呼叫 \texttt{setName("單位圓")}，
\texttt{main} 裡的 \texttt{c.setName(...)} 則編譯失敗。

\texttt{protected} 也可以用在成員變數上，例如把 \texttt{name} 直接標成
\texttt{protected}。 一般不建議這樣做：任何人寫一個衍生類別，就能任意改
\texttt{name}，要確認不變量就得檢查所有衍生類別。 成員變數維持
\texttt{private}，需要時開 \texttt{protected} 的成員函式。

\subsection{常見錯誤}

\begin{itemize}
\item \textbf{漏寫 \texttt{public}}：寫成 \texttt{class Circle : Shape}，
\texttt{class} 預設的繼承方式是 \texttt{private}，\texttt{Shape} 的成員在
外面全都不能用了。 \texttt{c.getName()} 報兩行：
\begin{code}{text}{編譯器的錯誤訊息}
error: 'const std::string& Shape::getName() const' is inaccessible within this context
error: 'Shape' is not an accessible base of 'Circle'
\end{code}
\texttt{private} 繼承很少用到。 本書所有的繼承都是 \texttt{public}，記得
每次都寫出來。
\item \textbf{衍生類別直接用基底的 \texttt{private} 成員}：在
\texttt{Circle} 裡寫 \texttt{void rename() \{ name = "圓形"; \}}，報
\texttt{'std::string Shape::name' is private within this context}。 改用
\texttt{Shape} 提供的函式。
\item \textbf{沒有呼叫基底的建構子}：把 \texttt{Circle} 的初始化列表寫成
\texttt{: r(r)}，編譯器會改用 \texttt{Shape} 的預設建構子，但
\texttt{Shape} 沒有，報 \texttt{no matching function for call to
'Shape::Shape()'}。 寫成 \texttt{: name("圓"), r(r)} 也不行，衍生類別的
初始化列表只能寫基底類別和自己的成員，報
\texttt{class 'Circle' does not have any field named 'name'}。
\end{itemize}

\section{建構與解構的順序}
\label{sec:inherit-order}

第~4~章看過，區塊裡的物件依宣告順序建立，以相反順序銷毀；第~5~章看過，成
員在建構子本體執行之前就建立好了。 有了繼承，衍生類別的物件裡有三個部分：
基底類別那一份、自己的成員、建構子本體。 下面的 \texttt{Tracer} 跟第~4~章
的一樣，建立和銷毀時各印一行：

\begin{example}[label=ex:inherit-order]{cpp}{(inherit\_order.cpp) 基底、成員、本體的順序}
#include <iostream>
#include <string>
using namespace std;

struct Tracer {
    string name;
    Tracer(const string &n) : name(n) { cout << "建立 " << name << '\n'; }
    ~Tracer() { cout << "銷毀 " << name << '\n'; }
};

class Base {
public:
    Base() : part("Base 的成員") { cout << "Base 建構子本體" << '\n'; }
    ~Base() { cout << "Base 解構子本體" << '\n'; }

private:
    Tracer part;
};

class Derived : public Base {
public:
    Derived() : part("Derived 的成員") { cout << "Derived 建構子本體" << '\n'; }
    ~Derived() { cout << "Derived 解構子本體" << '\n'; }

private:
    Tracer part;
};

int main() {
    {
        Derived d;
        cout << "--- d 建好了 ---" << '\n';
    }
    cout << "main 結束" << '\n';
    return 0;
}
\end{example}

\iosample{}{
建立 Base 的成員 \\
Base 建構子本體 \\
建立 Derived 的成員 \\
Derived 建構子本體 \\
--- d 建好了 --- \\
Derived 解構子本體 \\
銷毀 Derived 的成員 \\
Base 解構子本體 \\
銷毀 Base 的成員 \\
main 結束 \\
}

建立 \texttt{d} 時，先建立基底類別那一份（\texttt{Base} 自己也是先成員、
後本體），再依宣告順序建立 \texttt{Derived} 自己的成員，最後執行
\texttt{Derived} 建構子的本體。

銷毀時剛好反過來：先執行 \texttt{Derived} 解構子的本體，銷毀
\texttt{Derived} 的成員，最後才輪到 \texttt{Base} 那一份。 道理跟第~4~章一
樣：後建立的可能用到先建立的。 \texttt{Derived} 的建構子本體執行時，基底
和成員都已經建好，可以放心使用；\texttt{Derived} 的解構子本體執行時，它們
也都還在。

衍生類別的解構子不用、也不能自己呼叫基底的解構子，編譯器會在最後自動呼叫。
建構子則要在初始化列表選一個基底的建構子（上一節的 \texttt{Shape("圓")}），
沒寫的話就用基底的預設建構子；範例裡的 \texttt{Base} 只有預設建構子，所以
不用寫。

\subsection{常見錯誤}

\begin{itemize}
\item \textbf{以為初始化列表的順序決定建立順序}：把 \texttt{Derived} 的初
始化列表寫成 \texttt{: part("Derived 的成員"), Base()}，執行結果完全相同，
基底類別一定最先建立。 \texttt{-Wall} 會警告 \texttt{'Derived::part' will
be initialized after}，下一行是 \texttt{base 'Base'}。 跟第~5~章的成員
順序一樣，初始化列表照「基底、成員宣告順序」來排。
\item \textbf{在解構子裡自己呼叫基底的解構子}：在 \texttt{\textasciitilde
Derived()} 裡寫 \texttt{Base::\textasciitilde Base();}，\texttt{Base} 的解
構子會執行兩次。 \texttt{Base} 管著資源的話，就是第~5~章的重複釋放。
\end{itemize}

\section{虛擬函式}
\label{sec:virtual}

範例~\ref{ex:derived} 裡，透過 \texttt{Shape} 的參考不能呼叫
\texttt{area}，因為 \texttt{Shape} 沒有這個函式。 我們想要的是 C 回顧框
裡函式指標的效果：寫一個只認得 \texttt{Shape} 的函式，傳進圓就用圓的公式，
傳進矩形就用矩形的公式。

做法是在基底類別宣告 \texttt{area}，前面加上 \texttt{virtual}，成為
\textbf{虛擬函式}（virtual function）。 衍生類別寫一個同名、參數和
\texttt{const} 都一樣的函式，叫\textbf{覆寫}（override）。 下面的範例另外
放了一個沒有 \texttt{virtual} 的 \texttt{hello}，用來對照：

\begin{example}[label=ex:virtual]{cpp}{(virtual.cpp) 虛擬函式與一般函式}
#include <iostream>
#include <string>
using namespace std;

class Shape {
public:
    Shape(const string &name) : name(name) {}
    const string &getName() const { return name; }
    virtual double area() const { return 0; }                 // 虛擬函式
    void hello() const { cout << "Shape 的 hello" << '\n'; }  // 一般函式

private:
    string name;
};

class Circle : public Shape {
public:
    Circle(double r) : Shape("圓"), r(r) {}
    double area() const override { return 3.14159 * r * r; }
    void hello() const { cout << "Circle 的 hello" << '\n'; }

private:
    double r;
};

class Rect : public Shape {
public:
    Rect(double w, double h) : Shape("矩形"), w(w), h(h) {}
    double area() const override { return w * h; }

private:
    double w, h;
};

void report(const Shape &s) {
    cout << s.getName() << ' ' << s.area() << '\n';
}

int main() {
    Circle c(1);
    Rect r(3, 4);
    report(c);
    report(r);

    Shape *p = &c;
    p->hello();                // 一般函式：看 p 的型別
    c.hello();
    p = &r;
    cout << p->area() << '\n'; // 虛擬函式：看 p 指著的物件
    return 0;
}
\end{example}

\iosample{}{
圓 3.14159 \\
矩形 12 \\
Shape 的 hello \\
Circle 的 hello \\
12 \\
}

\texttt{report} 只知道參數是 \texttt{Shape}，\texttt{s.area()} 卻分別呼叫
了 \texttt{Circle} 和 \texttt{Rect} 的版本。 這裡要分兩種型別：

\begin{itemize}
\item \textbf{靜態型別}：變數宣告的型別，編譯時就知道。 \texttt{s} 的靜態
型別是 \texttt{const Shape \&}，\texttt{p} 是 \texttt{Shape *}。
\item \textbf{動態型別}：參考綁著的、指標指著的物件實際的型別，執行時才知
道。 第一次呼叫 \texttt{report} 時 \texttt{s} 的動態型別是
\texttt{Circle}，第二次是 \texttt{Rect}。
\end{itemize}

一般函式看靜態型別，所以 \texttt{p->hello()} 呼叫 \texttt{Shape} 的
\texttt{hello}，就算 \texttt{p} 指著 \texttt{Circle}。 虛擬函式透過參考或
指標呼叫時看動態型別，執行時才決定呼叫哪一個，這叫\textbf{動態繫結}
（dynamic binding）。 同一行 \texttt{s.area()} 依物件不同而有不同的行為，
就是多型。

\texttt{Circle} 的 \texttt{hello} 跟 \texttt{Shape} 的同名，但
\texttt{Shape} 的不是虛擬函式，所以這只是另一個函式，把基底的
\texttt{hello} 遮住了。 \texttt{c.hello()} 用的是 \texttt{Circle} 的，透過
\texttt{Shape *} 用的是 \texttt{Shape} 的，同一個物件兩種行為，容易搞混。
不要在衍生類別重新定義基底的一般函式。

\subsection{\texttt{override} 與 \texttt{final}}

衍生類別的 \texttt{area} 後面寫了 \texttt{override}。 它不改變程式的行
為，只是告訴編譯器「這個函式要覆寫基底的虛擬函式」。 如果基底沒有一樣的
虛擬函式可以覆寫，編譯器就報錯。 例如 \texttt{Circle} 的 \texttt{area} 漏
了 \texttt{const}：

\begin{code}{text}{編譯器的錯誤訊息}
error: 'double Circle::area()' marked 'override', but does not override
\end{code}

沒有 \texttt{const} 的 \texttt{area()} 跟有 \texttt{const} 的是兩個不同的
函式，沒寫 \texttt{override} 的話就會編譯成功，結果是 \texttt{report(c)}
印出 \texttt{圓 0}（見本節的常見錯誤）。 覆寫虛擬函式時一律寫
\texttt{override}。 基底類別裡的那一個寫 \texttt{virtual}；衍生類別裡寫了
\texttt{override} 就不用再寫 \texttt{virtual}，覆寫的函式自動也是虛擬的。

\texttt{final} 禁止再往下覆寫或繼承：函式後面寫 \texttt{final}，衍生類別
就不能再覆寫它；寫成 \texttt{class Circle final : public Shape}，就不能再
有類別繼承 \texttt{Circle}，否則報 \texttt{cannot derive from 'final' base
'Circle'}。 平常不常用。

虛擬函式背後的做法跟 C 回顧框的函式指標差不多：每個有虛擬函式的類別有一
張函式指標表，每個物件裡藏著一個指向自己類別那張表的指標，呼叫
\texttt{s.area()} 時先找到表，再從表裡拿出 \texttt{area} 的位址。 代價是
每個物件多一個指標（本書測試的環境下，只有一個 \texttt{string} 成員的類
別，加了虛擬函式後 \texttt{sizeof} 從 $32$ 變成 $40$），每次呼叫多查一次
表，一般的程式不用在意。 表由編譯器填，不會填錯，也不用強制轉型。

\subsection{常見錯誤}

\begin{itemize}
\item \textbf{忘了 \texttt{virtual}}：基底的 \texttt{area} 沒有寫
\texttt{virtual}，衍生類別也沒寫 \texttt{override}，程式照樣編譯成功。
\texttt{report} 和 \texttt{p->area()} 都用 \texttt{Shape} 的版本，三行都印
出 $0$。 衍生類別有寫 \texttt{override} 的話，編譯器會報上面的
\texttt{marked 'override', but does not override}，馬上就發現了。
\item \textbf{參數或 \texttt{const} 不一樣}：\texttt{Circle} 寫成
\texttt{double area()}，沒有 \texttt{const}，也沒有 \texttt{override}。 這
是一個新的函式，沒有覆寫任何東西，\texttt{report(c)} 印出 \texttt{圓 0}。
\texttt{g++ -Wall} 會警告 \texttt{'virtual double Shape::area() const' was
hidden [-Woverloaded-virtual=]}，但警告很容易被忽略，寫上
\texttt{override} 讓它變成錯誤。
\end{itemize}

\section{抽象類別與虛擬解構子}
\label{sec:abstract}

範例~\ref{ex:virtual} 的 \texttt{Shape::area} 回傳 $0$，其實是一個沒有意
義的值：「形狀」本身沒有面積，只有具體的圓或矩形才有。 這種基底類別無法
提供的函式，可以宣告成\textbf{純虛擬函式}（pure virtual function），在宣
告後面寫 \texttt{= 0}，不寫本體。

有純虛擬函式的類別叫\textbf{抽象類別}（abstract class），不能直接建立物
件，只能拿來當基底類別。 衍生類別要覆寫所有純虛擬函式，才能建立物件：

\begin{example}[label=ex:abstract]{cpp}{(abstract.cpp) 抽象類別 Shape}
#include <iostream>
#include <string>
using namespace std;

class Shape {                                   // 抽象類別
public:
    Shape(const string &name) : name(name) {}
    virtual ~Shape() = default;                 // 虛擬解構子，下面說明
    const string &getName() const { return name; }
    virtual double area() const = 0;            // 純虛擬函式
    virtual double perimeter() const = 0;

private:
    string name;
};

class Circle : public Shape {
public:
    Circle(double r) : Shape("圓"), r(r) {}
    double area() const override { return 3.14159 * r * r; }
    double perimeter() const override { return 2 * 3.14159 * r; }

private:
    double r;
};

class Rect : public Shape {
public:
    Rect(double w, double h) : Shape("矩形"), w(w), h(h) {}
    double area() const override { return w * h; }
    double perimeter() const override { return 2 * (w + h); }

private:
    double w, h;
};

void report(const Shape &s) {
    cout << s.getName() << " 面積 " << s.area() << " 周長 " << s.perimeter() << '\n';
}

int main() {
    Circle c(1);
    Rect r(3, 4);
    report(c);
    report(r);
    // Shape s("形狀");                         // 編譯失敗：Shape 是抽象類別
    return 0;
}
\end{example}

\iosample{}{
圓 面積 3.14159 周長 6.28318 \\
矩形 面積 12 周長 14 \\
}

\begin{itemize}
\item \texttt{Shape} 規定了「每種形狀都要能算面積和周長」，但不說怎麼算。
\texttt{report} 只依賴這個規定，之後多寫一個三角形，\texttt{report} 一行都
不用改。 C 回顧框的 \texttt{switch} 版本則要每個函式都補一個
\texttt{case}。
\item 拿掉最後一行的註解，\texttt{Shape s("形狀");} 編譯失敗，訊息還列出
了是哪些函式讓它變成抽象的：
\begin{code}{text}{編譯器的錯誤訊息}
error: cannot declare variable 's' to be of abstract type 'Shape'
note:   because the following virtual functions are pure within 'Shape':
note:     'virtual double Shape::area() const'
note:     'virtual double Shape::perimeter() const'
\end{code}
\end{itemize}

只有純虛擬函式（加上虛擬解構子）、沒有成員變數的抽象類別，常叫做
\textbf{介面}（interface），它只規定衍生類別要提供哪些函式。

\subsection{虛擬解構子}

\texttt{Shape} 裡有一行 \texttt{virtual \textasciitilde Shape() = default;}。
\texttt{= default} 是第~6~章說的「用編譯器產生的版本」，這裡要的是前面那個
\texttt{virtual}。 原因是物件常常透過基底類別的指標銷毀，例如
\texttt{unique\_ptr<Shape>} 指著一個 \texttt{Circle}。 \texttt{unique\_ptr}
銷毀時對它的指標做 \texttt{delete}，這個指標的型別是 \texttt{Shape *}。
解構子不是虛擬函式的話，\texttt{delete} 就看靜態型別，只呼叫
\texttt{Shape} 的解構子。 下面的 \texttt{Derived} 有一個 \texttt{vector}
成員，基底 \texttt{Base} 的解構子沒有 \texttt{virtual}：

\begin{example}[label=ex:vdtor]{cpp}{(vdtor.cpp) 基底的解構子不是 virtual（錯誤示範）}
#include <iostream>
#include <memory>
#include <vector>
using namespace std;

class Base {
public:
    ~Base() { cout << "~Base" << '\n'; }        // 沒有 virtual
    virtual void work() const { cout << "Base::work" << '\n'; }
};

class Derived : public Base {
public:
    Derived() : buf(1000, 7) {}
    ~Derived() { cout << "~Derived" << '\n'; }
    void work() const override { cout << "Derived::work " << buf.size() << '\n'; }

private:
    vector<int> buf;
};

int main() {
    unique_ptr<Base> p = make_unique<Derived>();
    p->work();
    return 0;
}                                               // p 銷毀：delete 一個 Base *
\end{example}

\iosample{}{
Derived::work 1000 \\
\textasciitilde Base \\
}

\texttt{g++ -Wall -Wextra} 沒有任何警告，執行也沒有當掉，但
\texttt{\textasciitilde Derived} 沒有印出來：\texttt{Derived} 的解構子沒有
執行，\texttt{buf} 也就沒有銷毀。 用 ASan 跑：

\begin{code}{text}{用 ASan 編譯後執行（節錄）}
Derived::work 1000
~Base
==2041472==ERROR: AddressSanitizer: new-delete-type-mismatch on 0x503000000040 in thread T0:
  object passed to delete has wrong type:
  size of the allocated type:   32 bytes;
  size of the deallocated type: 8 bytes.
    #3 0x5cb786796532 in main vdtor.cpp:26

allocated by thread T0 here:
    #2 0x5cb786796492 in main vdtor.cpp:23
\end{code}

\texttt{new-delete-type-mismatch} 是 \texttt{delete} 的型別跟
\texttt{new} 的不一樣：第 $23$ 行配置的是 $32$ 個位元組的 \texttt{Derived}，
第 $26$ 行（\texttt{p} 銷毀）卻當成 $8$ 個位元組的 \texttt{Base} 釋放。
照標準，透過基底指標 \texttt{delete} 衍生類別的物件，而基底的解構子不是虛
擬的，是未定義行為。 這個例子實際的後果是 \texttt{buf} 那 $4000$ 個位元組
洩漏了（關掉這項檢查再跑，LeakSanitizer 報的就是
\texttt{Direct leak of 4000 byte(s)}）。

把第 $8$ 行改成 \texttt{virtual \textasciitilde Base() \{ ... \}}，
\texttt{delete} 就看動態型別，先呼叫 \texttt{\textasciitilde Derived}，再
呼叫 \texttt{\textasciitilde Base}（上一節的順序），輸出多一行
\texttt{\textasciitilde Derived}，ASan 也沒有報告。 規則很簡單：\textbf{有虛
擬函式的基底類別，解構子一定要是 \texttt{virtual}}。 不需要做事的話，寫
\texttt{virtual \textasciitilde Shape() = default;} 就好。

第~6~章說過，自己宣告了解構子，編譯器就不產生移動。 虛擬解構子也是，但通
常沒有關係：多型的物件透過指標使用，不會整個複製或移動（原因見下一節）。

\subsection{常見錯誤}

\begin{itemize}
\item \textbf{漏了一個純虛擬函式}：\texttt{Rect} 沒有寫 \texttt{perimeter}，
它就還是抽象類別，\texttt{Rect r(3, 4);} 報 \texttt{cannot declare variable
'r' to be of abstract type 'Rect'}，\texttt{note} 列出還沒實作的
\texttt{'virtual double Shape::perimeter() const'}。
\item \textbf{以為 \texttt{unique\_ptr} 會處理好}：上面的範例用的就是
\texttt{unique\_ptr}，照樣只呼叫 \texttt{\textasciitilde Base}。
\texttt{unique\_ptr} 保證會 \texttt{delete}，但 \texttt{delete} 呼叫哪個解
構子，要看基底的解構子是不是虛擬的。 如果是自己寫 \texttt{Base *p = new
Derived; delete p;}，\texttt{-Wall} 會警告 \texttt{deleting object of
polymorphic class type 'Base' which has non-virtual destructor might cause
undefined behavior}；包在 \texttt{unique\_ptr} 裡反而看不到這個警告。
\item \textbf{在建構子或解構子裡呼叫虛擬函式}：\texttt{Base} 的建構子執行
時，\texttt{Derived} 的部分還沒建立（第~\ref{sec:inherit-order}~節），所以
這時呼叫虛擬函式，用的是 \texttt{Base} 的版本，不會呼叫到
\texttt{Derived} 的。 解構子也一樣。 呼叫的是純虛擬函式的話，結果是未定義
行為。 不要在建構子和解構子裡依賴虛擬函式（本章練習第~1~題）。
\end{itemize}

\section{用 \texttt{vector<unique\_ptr<Base>>} 存放多型物件}
\label{sec:poly-container}

C 回顧框的 \texttt{switch} 版本可以把所有形狀放進一個
\texttt{struct Shape} 陣列。 用繼承之後，圓和矩形是不同的型別，大小也不一
樣，不能直接放進同一個 \texttt{vector}。 能放進去的是指向它們的指標：
\texttt{vector<unique\_ptr<Shape>>}。 每個元素都是 \texttt{unique\_ptr<Shape>}，
指著的物件可以是任何一種形狀。

下面的程式讀入一串形狀，算出總面積和最大的那一個。 \texttt{Shape}、
\texttt{Circle}、\texttt{Rect} 跟範例~\ref{ex:abstract} 相同：

\begin{example}[label=ex:poly-vector]{cpp}{(shapes.cpp) 把不同的形狀放進同一個 vector}
#include <iostream>
#include <memory>
#include <string>
#include <vector>
using namespace std;

// Shape、Circle、Rect 同範例 abstract.cpp，這裡省略

unique_ptr<Shape> makeShape(const string &kind) {
    if (kind == "circle") {
        double r;
        cin >> r;
        return make_unique<Circle>(r);
    }
    if (kind == "rect") {
        double w, h;
        cin >> w >> h;
        return make_unique<Rect>(w, h);
    }
    return nullptr;
}

int main() {
    vector<unique_ptr<Shape>> shapes;
    string kind;
    while (cin >> kind) {
        unique_ptr<Shape> s = makeShape(kind);
        if (s == nullptr) {
            break;                              // 不認得的種類就停止讀取
        }
        shapes.push_back(std::move(s));
    }

    double total = 0;
    const Shape *largest = nullptr;             // 只用來看，不擁有
    for (const unique_ptr<Shape> &s : shapes) {
        cout << s->getName() << ' ' << s->area() << '\n';
        total += s->area();
        if (largest == nullptr || s->area() > largest->area()) {
            largest = s.get();
        }
    }
    cout << "總面積 " << total << '\n';
    if (largest != nullptr) {
        cout << "最大的是" << largest->getName() << '\n';
    }
    return 0;
}                                               // shapes 銷毀，每個形狀跟著銷毀
\end{example}

\iosample{
circle 1 \\
rect 3 4 \\
circle 2 \\
}{
圓 3.14159 \\
矩形 12 \\
圓 12.5664 \\
總面積 27.708 \\
最大的是圓 \\
}

\begin{itemize}
\item \texttt{makeShape} 依讀到的種類建立不同的物件，回傳型別都是
\texttt{unique\_ptr<Shape>}。 \texttt{make\_unique<Circle>(r)} 建立的是
\texttt{unique\_ptr<Circle>}，可以自動轉成 \texttt{unique\_ptr<Shape>}，就
像 \texttt{Circle *} 可以轉成 \texttt{Shape *}。 這種「依資料決定建立哪個
衍生類別」的函式，常叫做\textbf{工廠函式}（factory function）。
\item \texttt{unique\_ptr} 不能複製，所以放進 \texttt{vector} 時要
\texttt{std::move}（第~4~章），走訪時用 \texttt{const} 參考。
\texttt{s->area()} 透過指標呼叫虛擬函式，依每個元素實際的型別呼叫。
\item \texttt{largest} 是一般指標，用 \texttt{s.get()} 取得。 它只是記住
「最大的是哪一個」，不擁有物件，也不負責 \texttt{delete}；物件仍然由
\texttt{vector} 裡的 \texttt{unique\_ptr} 擁有。 \texttt{shapes} 還在，
\texttt{largest} 就有效。
\item \texttt{shapes} 銷毀時，每個 \texttt{unique\_ptr} 透過
\texttt{Shape *} 刪除自己的物件；\texttt{Shape} 有虛擬解構子，所以都正確
銷毀，用 ASan 跑沒有報告。
\end{itemize}

多一種形狀，只要寫一個新的衍生類別，再在 \texttt{makeShape} 加一個
\texttt{if}；\texttt{main} 裡的迴圈不用改。

\subsection{物件切割}

如果不用指標，直接寫 \texttt{vector<Shape>} 會怎樣？ \texttt{Shape} 是抽
象類別的話，編譯器不讓你這樣做。 換成範例~\ref{ex:virtual} 那個不是抽象的
\texttt{Shape}（\texttt{area} 回傳 $0$），把它的 \texttt{main} 換成下面這
段（記得 \texttt{\#include <vector>}）：

\begin{code}{cpp}{(slicing.cpp) 用傳值的方式存 Circle 和 Rect（錯誤示範）}
int main() {
    Circle c(1);
    Shape s = c;               // 只複製 c 裡 Shape 的那一部分
    cout << s.getName() << ' ' << s.area() << '\n';

    vector<Shape> v;
    v.push_back(c);
    v.push_back(Rect(3, 4));
    for (const Shape &x : v) {
        report(x);
    }
    return 0;
}
\end{code}

編譯沒有警告，執行結果是 \texttt{圓 0}、\texttt{圓 0}、\texttt{矩形 0} 三
行，面積全部變成 $0$。 \texttt{Shape s = c;} 建立的是一個
\texttt{Shape} 物件，用 \texttt{Shape} 的複製建構子，只複製 \texttt{c} 裡
\texttt{Shape} 的那一部分（\texttt{name}），\texttt{r} 被切掉了。 \texttt{s}
從頭到尾就是一個 \texttt{Shape}，呼叫的當然是 \texttt{Shape::area}。 這叫
\textbf{物件切割}（object slicing）。 \texttt{vector<Shape>} 的每個元素也
都是 \texttt{Shape}，\texttt{push\_back(c)} 一樣只複製 \texttt{Shape} 的部
分。 傳值的參數 \texttt{void report(Shape s)} 也會切割。

要保留動態型別，就不能複製出新的物件，只能透過參考或指標使用原本的物件：
參數用 \texttt{const Shape \&}，容器用 \texttt{vector<unique\_ptr<Shape>>}。
這也是基底類別最好寫成抽象類別的一個理由：\texttt{Shape} 是抽象的時候，
上面的 \texttt{Shape s = c;} 直接編譯失敗（\texttt{cannot declare variable
's' to be of abstract type 'Shape'}），傳值的參數也是
（\texttt{cannot declare parameter 's' to be of abstract type 'Shape'}）。

\subsection{\texttt{dynamic\_cast}}

偶爾需要知道一個 \texttt{Shape} 實際上是不是 \texttt{Circle}，可以寫
\texttt{const Circle *c = dynamic\_cast<const Circle *>(s.get());}。
\texttt{s} 指著的是 \texttt{Circle} 時，\texttt{c} 指向它，可以用
\texttt{Circle} 自己的成員；不是的話，\texttt{c} 是 \texttt{nullptr}。
\texttt{dynamic\_cast} 只能用在有虛擬函式的類別上。 程式裡常常需要
\texttt{dynamic\_cast}，通常表示基底類別少了一個虛擬函式：與其問「你是不是
圓」再分別處理，不如在 \texttt{Shape} 加一個虛擬函式，讓每種形狀自己處理。

\subsection{常見錯誤}

\begin{itemize}
\item \textbf{\texttt{vector<Shape *>} 加上 \texttt{new}}：用一般指標也能做到
多型，但每個元素都要記得 \texttt{delete}，中途 \texttt{return} 或出例外就洩
漏，這是第~4~章的問題。 用 \texttt{unique\_ptr}。
\item \textbf{把 \texttt{s.get()} 交給另一個 \texttt{unique\_ptr}}：
\texttt{unique\_ptr<Shape> q(shapes[0].get());} 讓兩個
\texttt{unique\_ptr} 都以為自己擁有同一個物件，最後 \texttt{delete} 兩次。
\texttt{get()} 取出來的指標只能拿來看。
\end{itemize}

\section{什麼時候不要用繼承}
\label{sec:composition}

繼承容易被拿來「重複使用程式碼」：某個類別已經有我要的函式，繼承它就全部
有了。 下面是一個常見的例子，想要一個堆疊（只能從頂端放入和取出），於是
直接繼承 \texttt{vector<int>}。 對照的 \texttt{Stack} 則是把
\texttt{vector} 放成成員，這叫\textbf{組合}（composition）：

\begin{example}[label=ex:stack]{cpp}{(stack.cpp) 繼承 vector 與把 vector 當成員}
#include <iostream>
#include <vector>
using namespace std;

class BadStack : public vector<int> {   // 繼承 vector，不建議
public:
    void push(int x) { push_back(x); }
    void pop() { pop_back(); }
    int top() const { return back(); }
};

class Stack {                           // vector 當成員
public:
    void push(int x) { data.push_back(x); }
    void pop() { data.pop_back(); }
    int top() const { return data.back(); }
    int size() const { return (int)data.size(); }

private:
    vector<int> data;
};

int main() {
    BadStack b;
    b.push(1);
    b.push(2);
    b[0] = 99;                  // 連堆疊底部都改得到
    b.push_back(3);
    cout << b.top() << ' ' << b[0] << ' ' << b.size() << '\n';

    Stack s;
    s.push(1);
    s.push(2);
    s.push(3);
    s.pop();
    // s[0] = 99;               // 編譯失敗
    cout << s.top() << ' ' << s.size() << '\n';
    return 0;
}
\end{example}

\iosample{}{
3 99 3 \\
2 2 \\
}

\begin{itemize}
\item \texttt{BadStack} 是一種 \texttt{vector}，所以 \texttt{vector} 的
\texttt{public} 成員函式它全都有：\texttt{b[0] = 99} 改了堆疊最底下的元素，
\texttt{push\_back} 繞過了 \texttt{push}。 「只能從頂端存取」這條規則，類
別本身擋不住。
\item \texttt{Stack} 有一個 \texttt{vector}（has-a），外面只看得到
\texttt{Stack} 自己開放的四個函式。 \texttt{s[0] = 99} 報 \texttt{no match
for 'operator[]' (operand types are 'Stack' and 'int')}。 之後想把
\texttt{vector} 換成別的容器，用 \texttt{Stack} 的程式也不用改。
\item 標準函式庫的容器沒有虛擬解構子，本來就不是設計來當基底類別的。
\end{itemize}

判斷的方法是第~\ref{sec:derived}~節那句話：「\texttt{A} 是一種
\texttt{B}」說得通，而且任何需要 \texttt{B} 的地方交給它一個 \texttt{A} 都
不會出問題，才用 \texttt{public} 繼承。 只是「\texttt{A} 用到
\texttt{B}」或「\texttt{A} 有一個 \texttt{B}」，就把 \texttt{B} 放成成員。
堆疊不是一種 \texttt{vector}，任何 \texttt{vector} 能做的事堆疊都不該做。

\subsection{型別標記加 \texttt{switch} 有時候比較好}

本章開頭 C 回顧框的第一種寫法，在 C++ 裡也不算錯。 型別標記改用
\texttt{enum class Kind \{ Circle, Rect, Triangle \};}，\texttt{switch} 裡
沒有 \texttt{default}，只寫了 \texttt{Circle} 和 \texttt{Rect} 兩個
\texttt{case} 的話，\texttt{g++ -Wall} 會警告 \texttt{enumeration value
'Triangle' not handled in switch [-Wswitch]}。 多一種形狀時，編譯器會指出
每一個要補的 \texttt{switch}。 兩種寫法各有適合的情況：

\begin{itemize}
\item 種類會一直增加，對每個物件做的操作比較固定（例如外掛、遊戲裡的各種
角色）：用虛擬函式。 多一種，只要多寫一個類別。
\item 種類固定只有幾種，操作卻會一直增加（例如一個小計算機的運算子種類、
一份資料只有三種紀錄格式）：用型別標記加 \texttt{switch}。 多一種操作，只
要多寫一個函式，不用改每個類別。
\end{itemize}

整體的建議是：先想清楚是不是真的需要多型，需要的話保持繼承階層又淺又
小，一個抽象的基底類別加上幾個衍生類別，就能處理大部分情況。 只是想重複
使用程式碼，用組合。

\subsection{常見錯誤}

\begin{itemize}
\item \textbf{為了少寫幾行而繼承}：\texttt{class Student : public
vector<int>}，因為學生「有一些分數」。 學生不是一種 \texttt{vector}，分
數應該是成員。
\item \textbf{很深的繼承階層}：\texttt{Shape} 下面有 \texttt{Polygon}，
下面有 \texttt{Quadrilateral}，下面有 \texttt{Rect}，再下面有
\texttt{Square}。 讀一個函式要往上翻好幾層才知道它從哪裡來，改基底的一行
可能影響所有衍生類別。 層數多到要畫圖才看得懂時，通常可以用組合拆開。
\end{itemize}

\section{本章練習}
\label{sec:ch7-exercises}

\begin{drills}[label=drills:ch7]
\item 預測下面程式的輸出。
\begin{code}{cpp}{練習：建構、虛擬函式與切割}
#include <iostream>
using namespace std;

struct A {
    A() { cout << "A "; hi(); }
    virtual ~A() { cout << "~A "; }
    virtual void hi() const { cout << "A.hi "; }
    void name() const { cout << "A.name "; }
};

struct B : A {
    B() { cout << "B "; hi(); }
    ~B() override { cout << "~B "; }
    void hi() const override { cout << "B.hi "; }
    void name() const { cout << "B.name "; }
};

void call(const A &a) {
    a.hi();
    a.name();
}

int main() {
    B b;
    cout << '\n';
    call(b);
    b.name();
    cout << '\n';
    A a = b;
    a.hi();
    cout << '\n';
    return 0;
}
\end{code}
\item 下面的程式編譯失敗，\texttt{g++} 報了四個錯誤。 找出原因並改正，讓
它輸出 \texttt{Lucky: 汪} 和 \texttt{Momo: 喵} 兩行，而且 \texttt{Cat} 仍
然可以改名。 改到編譯成功之後，還有一個問題是 \texttt{g++} 和 ASan 都不會
告訴你的，也一起改掉。
\begin{code}{cpp}{練習：找錯}
#include <iostream>
#include <memory>
#include <string>
#include <vector>
using namespace std;

class Animal {
public:
    Animal(const string &name) : name(name) {}
    virtual string sound() const = 0;
    void speak() const { cout << name << ": " << sound() << '\n'; }

private:
    string name;
};

class Dog : Animal {
public:
    Dog(const string &name) : Animal(name) {}
    string sound() override { return "汪"; }
};

class Cat : public Animal {
public:
    Cat(const string &name) : Animal(name) {}
    string sound() const override { return "喵"; }
    void rename(const string &n) { name = n; }
};

int main() {
    vector<unique_ptr<Animal>> pets;
    pets.push_back(make_unique<Dog>("Lucky"));
    auto cat = make_unique<Cat>("Mimi");
    cat->rename("Momo");
    pets.push_back(std::move(cat));
    for (const unique_ptr<Animal> &p : pets) {
        p->speak();
    }
    return 0;
}
\end{code}
\item 停車場的收費方式：汽車每小時 $40$ 元；機車每小時 $20$ 元，但最多收
$100$ 元。 寫一個抽象類別 \texttt{Vehicle}，建構子收停車的時數，提供
\texttt{getHours()}，以及純虛擬函式 \texttt{kind()}（回傳 \texttt{"汽車"}
或 \texttt{"機車"}）和 \texttt{fee()}；再寫兩個衍生類別 \texttt{Car} 和
\texttt{Motorcycle}。 \texttt{main} 每行讀入種類（\texttt{car} 或
\texttt{moto}）和時數直到輸入結束，不認得的種類略過；全部讀完後，依輸入順
序印出每一輛的種類、時數和費用，最後印出合計。 輸入
\begin{code}{text}{練習 3 的輸入}
car 3
moto 2
moto 8
bus 5
car 1
\end{code}
應該輸出 \texttt{汽車 3 小時 120}、\texttt{機車 2 小時 40}、
\texttt{機車 8 小時 100}、\texttt{汽車 1 小時 40}、\texttt{合計 300} 五行。
\item 下面的程式用 \texttt{g++ -Wall -Wextra} 編譯沒有警告，輸出
\texttt{Ada 990} 和 \texttt{Alan 4990}。 \texttt{Alan} 的是儲蓄帳戶，月底
應該加 $1\%$ 利息、不扣手續費，變成 $5050$ 才對。 說明原因並改正。
\begin{code}{cpp}{練習：月底結算}
#include <iostream>
#include <string>
#include <vector>
using namespace std;

class Account {
public:
    Account(const string &owner, int balance) : owner(owner), balance(balance) {}
    virtual ~Account() = default;
    void deposit(int amount) { balance += amount; }
    int getBalance() const { return balance; }
    virtual void monthEnd() { balance -= 10; }   // 每月手續費 10 元
    void print() const { cout << owner << ' ' << balance << '\n'; }

private:
    string owner;
    int balance;
};

class Savings : public Account {
public:
    Savings(const string &owner, int balance) : Account(owner, balance) {}
    void monthEnd() override { deposit(getBalance() / 100); }   // 利息 1%，免手續費
};

int main() {
    vector<Account> accounts;
    accounts.push_back(Account("Ada", 1000));
    accounts.push_back(Savings("Alan", 5000));
    for (Account &a : accounts) {
        a.monthEnd();
    }
    for (const Account &a : accounts) {
        a.print();
    }
    return 0;
}
\end{code}
\item 下面每一組，判斷應該用 \texttt{public} 繼承、組合（放成成員），還是
型別標記加 \texttt{switch}，並說明理由。
\begin{enumerate}
\item \texttt{Car} 和 \texttt{Engine}（引擎）。
\item \texttt{Student} 和 \texttt{Person}，程式裡有些函式收
\texttt{const Person \&}，要能把學生傳進去。
\item 一個可以改寬和高的 \texttt{Rect}（有 \texttt{setWidth}、
\texttt{setHeight}），想再寫一個 \texttt{Square}。
\item \texttt{Queue}（佇列，只能從尾端放入、從前端取出）和
\texttt{vector<int>}。
\item 一個記帳程式的每筆紀錄只有「收入」和「支出」兩種，之後會一直加新的
報表功能（月報、分類統計、匯出）。
\end{enumerate}
\end{drills}

\addcontentsline{toc}{section}{練習解答}
\section*{練習解答}

\begin{enumerate}
\item 輸出四行：\texttt{A A.hi B B.hi }、\texttt{B.hi A.name B.name }、
\texttt{A.hi }、\texttt{\textasciitilde A \textasciitilde B \textasciitilde A }
（最後一行沒有換行）。 建立 \texttt{b} 時先建立 \texttt{A} 的部分，這時
\texttt{B} 的部分還不存在，所以 \texttt{A} 建構子裡的 \texttt{hi()} 呼叫的
是 \texttt{A.hi}；接著 \texttt{B} 的建構子印 \texttt{B}，這時呼叫
\texttt{hi()} 才是 \texttt{B.hi}。 \texttt{call(b)} 裡 \texttt{a.hi()} 是
虛擬函式，看動態型別，印 \texttt{B.hi}；\texttt{a.name()} 不是虛擬函式，
看靜態型別 \texttt{A}，印 \texttt{A.name}。 \texttt{b.name()} 的靜態型別
是 \texttt{B}，印 \texttt{B.name}。 \texttt{A a = b;} 是物件切割，用編譯器
產生的複製建構子（不印東西），\texttt{a} 就是一個 \texttt{A}，所以
\texttt{a.hi()} 印 \texttt{A.hi}。 \texttt{main} 結束時先銷毀後建立的
\texttt{a}（\texttt{\textasciitilde A}），再銷毀 \texttt{b}
（\texttt{\textasciitilde B} 然後 \texttt{\textasciitilde A}）。
\item 四個錯誤訊息來自三個問題：
(1) \texttt{class Dog : Animal} 少了 \texttt{public}，是 \texttt{private}
繼承，\texttt{unique\_ptr<Dog>} 不能轉成 \texttt{unique\_ptr<Animal>}。 這
個錯誤報在 \texttt{pets.push\_back(...)} 那一行，訊息是
\texttt{no matching function for call to ...push\_back(...)}，後面的
\texttt{note} 很長，看不太出原因。 改成 \texttt{class Dog : public Animal}。
(2) \texttt{Dog::sound} 少了 \texttt{const}，報 \texttt{marked 'override',
but does not override}。 寫了 \texttt{override}，這個錯誤才被抓到。 因為
沒有覆寫到，\texttt{Dog} 還是抽象類別，所以第四個錯誤
\texttt{invalid new-expression of abstract class type 'Dog'}（位置在
\texttt{unique\_ptr.h}，是 \texttt{make\_unique<Dog>} 裡的 \texttt{new}）
也是同一個原因。
(3) \texttt{Cat::rename} 直接改 \texttt{Animal} 的 \texttt{private} 成員
\texttt{name}。 在 \texttt{Animal} 加一個 \texttt{protected} 的
\texttt{void setName(const string \&n) \{ name = n; \}}，\texttt{rename}
改成呼叫 \texttt{setName(n)}。 不要把 \texttt{name} 改成
\texttt{protected} 或 \texttt{public}。
還有一個問題：\texttt{Animal} 有虛擬函式，物件透過
\texttt{unique\_ptr<Animal>} 銷毀，解構子卻不是虛擬的。 這個程式裡
\texttt{Dog}、\texttt{Cat} 沒有自己的成員，大小跟 \texttt{Animal} 一樣，所以
ASan 剛好沒有報告，但這仍然是未定義行為；之後替 \texttt{Cat} 加一個
\texttt{string} 成員，它就會洩漏。 在 \texttt{Animal} 加上
\texttt{virtual \textasciitilde Animal() = default;}。
\item 一種寫法：
\begin{code}{cpp}{練習 3 參考解答（類別部分）}
class Vehicle {
public:
    Vehicle(int hours) : hours(hours) {}
    virtual ~Vehicle() = default;
    int getHours() const { return hours; }
    virtual string kind() const = 0;
    virtual int fee() const = 0;

private:
    int hours;
};

class Car : public Vehicle {
public:
    Car(int hours) : Vehicle(hours) {}
    string kind() const override { return "汽車"; }
    int fee() const override { return 40 * getHours(); }
};

class Motorcycle : public Vehicle {
public:
    Motorcycle(int hours) : Vehicle(hours) {}
    string kind() const override { return "機車"; }
    int fee() const override {
        int f = 20 * getHours();
        return f > 100 ? 100 : f;
    }
};
\end{code}
時數是每輛車都有的，放在 \texttt{Vehicle}，\texttt{private} 加上 getter；
收費方式每種不同，寫成純虛擬函式。 \texttt{Vehicle} 有虛擬函式，又透過
\texttt{unique\_ptr<Vehicle>} 銷毀，所以要有虛擬解構子。 \texttt{main} 照範
例~\ref{ex:poly-vector} 的寫法：寫一個工廠函式 \texttt{makeVehicle(type,
hours)}，\texttt{"car"} 回傳 \texttt{make\_unique<Car>(hours)}，
\texttt{"moto"} 回傳 \texttt{make\_unique<Motorcycle>(hours)}，其他回傳
\texttt{nullptr}；\texttt{while (cin >> type >> hours)} 讀入，不是
\texttt{nullptr} 才 \texttt{push\_back(std::move(v))} 進
\texttt{vector<unique\_ptr<Vehicle>> lot}；最後走訪 \texttt{lot}，印出
\texttt{v->kind()}、\texttt{v->getHours()}、\texttt{v->fee()} 並累加合計。

\item \texttt{vector<Account>} 的元素都是 \texttt{Account}。
\texttt{push\_back(Savings("Alan", 5000))} 只把 \texttt{Savings} 裡
\texttt{Account} 的部分複製進去，這是物件切割，元素已經不是
\texttt{Savings} 了，所以 \texttt{monthEnd} 用的是 \texttt{Account} 的版本，
扣了 $10$ 元手續費。 \texttt{Account} 不是抽象類別，所以編譯器不會擋。
改法是元素改存指標：
\begin{code}{cpp}{練習 4：改用 vector<unique\_ptr<Account>>}
vector<unique_ptr<Account>> accounts;
accounts.push_back(make_unique<Account>("Ada", 1000));
accounts.push_back(make_unique<Savings>("Alan", 5000));
for (const unique_ptr<Account> &a : accounts) {
    a->monthEnd();
}
for (const unique_ptr<Account> &a : accounts) {
    a->print();
}
\end{code}
記得 \texttt{\#include <memory>}。 改完輸出 \texttt{Ada 990} 和
\texttt{Alan 5050}。 第一個迴圈的 \texttt{a} 是 \texttt{const} 參考，但
\texttt{const} 的是 \texttt{unique\_ptr} 本身（不能改指向誰），不是它指著
的 \texttt{Account}，所以還是可以呼叫會改餘額的 \texttt{monthEnd}。
\item
(a) 組合。 汽車有一個引擎，不是一種引擎，\texttt{Engine} 放成
\texttt{Car} 的成員。
(b) \texttt{public} 繼承。 學生是一種人，需要 \texttt{Person} 的地方交給它
學生也說得通，正好是繼承要表達的關係。
(c) 不要讓 \texttt{Square} 繼承這個 \texttt{Rect}。 正方形在數學上是一種矩
形，但這個 \texttt{Rect} 可以只改寬不改高：一個函式收 \texttt{Rect \&}，
呼叫 \texttt{setWidth(5)} 之後預期高度不變，傳進正方形就出問題（要嘛高度
跟著變，要嘛它不再是正方形）。 「任何需要 \texttt{Rect} 的地方都能交給它
一個 \texttt{Square}」不成立。 可以讓 \texttt{Square} 和 \texttt{Rect} 都
直接繼承 \texttt{Shape}，或者讓形狀建立後不能改大小。
(d) 組合，跟 \texttt{Stack} 一樣。 佇列不該讓人用 \texttt{[]} 改中間的元
素。
(e) 型別標記加 \texttt{switch}。 種類固定只有兩種，會增加的是操作；每多一
個報表就是多一個函式，用 \texttt{enum class} 的話，漏了哪一種
\texttt{-Wall} 也會提醒。
\end{enumerate}
